\begin{proposition}[5.1.10]
\label{IV.2.5.1.10-fr}
\begin{enumerate}
\item[\textup{(i)}] En todo preesquema $X$, toda parte localmente construible no
vacía $E$ contiene un punto $x$ tal que $\{x\}$ sea una parte localmente
cerrada de $X$ \textup{(o, lo que equivale a lo mismo, tal que $x$ esté \emph{aislado en}
$\overline{\{x\}}$)}.
\item[\textup{(ii)}] Sean $X$ un preesquema localmente noetheriano y $x$ un punto de
$X$ tal que $\{x\}$ sea localmente cerrado en $X$; entonces se tiene
$\dim(\overline{\{x\}})\leq1$; por consiguiente, todo punto $y\ne x$ de
$\overline{\{x\}}$ es cerrado en $X$.
\end{enumerate}
\end{proposition}

\begin{proof}
(i) El resultado es un caso particular del lema siguiente:
\end{proof}

\begin{lemma}[5.1.10.1]
\label{IV.2.5.1.10.1-fr}
Sea $X$ un espacio topológico con la propiedad siguiente: para toda parte
localmente cerrada $Z\ne\varnothing$ de $X$, existen una parte $Z'$ de $Z$,
localmente cerrada en $X$ \textup{(o \emph{en} $Z$, lo que equivale a lo mismo)}, y un punto
$x\in Z'$ cerrado \textup{en} $Z'$. Entonces toda parte localmente cerrada
$Z\ne\varnothing$ de $X$ contiene un punto $x$ aislado en
$\overline{\{x\}}$.
\end{lemma}

\begin{proof}
En efecto, sea $Z'\subset Z$ una parte localmente cerrada de $Z$ que contiene
un punto $x$ tal que $x$ sea cerrado \emph{en} $Z'$. Existe una vecindad abierta
$U$ de $x$ en $X$ tal que $Z'\cap U$ sea cerrado \emph{en} $U$, por lo que $x$ es
también cerrado en $U$; esto significa que $U\cap\overline{\{x\}}=\{x\}$, es
decir, que $x$ está aislado en $\overline{\{x\}}$.

El lema se aplica a todo espacio subyacente a un preesquema $X$, pues $Z$ es
entonces también el espacio subyacente a un preesquema (\textbf I, 5.2.1), y basta
tomar por $Z'$ un abierto afín \emph{en} $Z$, que es un espacio de Kolmogoroff
cuasicompacto y, por tanto, contiene un punto cerrado
$(\mathbf{0}_{\mathrm I},2.1.3)$.
\end{proof}

\oldpage[IV-2]{89}

\begin{proof}
\textup{(ii)} Sea $Z$ el subpreesquema reducido de $X$ que tiene a
$\overline{\{x\}}$ por espacio
subyacente; la hipótesis implica que $\{x\}$ es abierto \emph{en} $Z$;
por tanto, para todo $z\in Z$, el punto genérico $x$ de
$\operatorname{Spec}(\mathcal O_{Z,z})$ está \emph{aislado} en
$\operatorname{Spec}(\mathcal O_{Z,z})$; pero el anillo
$A=\mathcal O_{Z,z}$ es un anillo local noetheriano íntegro, y la hipótesis
implica que existe un $f\in A$ tal que $A_f$ sea un \emph{cuerpo}; se sabe
$(\mathbf{0},16.3.3)$ que esto implica $\dim(A)\leq1$. Como
$\dim(\mathcal O_{Z,z})\leq1$ para todo $z\in Z$, se tiene efectivamente
$\dim(Z)\leq1$.
\end{proof}

\begin{corollary}[5.1.11]
\label{IV.2.5.1.11-fr}
Si $X$ es un preesquema localmente noetheriano, toda parte cerrada no vacía de
$X$ contiene un punto cerrado.
\end{corollary}

\begin{proof}
En efecto, toda parte cerrada de $X$ es construible
$(\mathbf{0}_{\mathrm{III}},9.1.1\text{ y }9.1.5)$, y basta aplicar (5.1.10).
\end{proof}
